% MỞ ĐẦU (LỜI MỞ ĐẦU - Chương 0)

\chapter*{MỞ ĐẦU} % Tên của chương
\addcontentsline{toc}{chapter}{MỞ ĐẦU} % Thêm tên chương vào mục lục

\label{Chapter0} % Để trích dẫn chương này ở chỗ nào đó trong bài, hãy sử dụng lệnh \ref{Chapter0} 

%----------------------------------------------------------------------------------------

Trong thời đại số, lượng dữ liệu văn bản khổng lồ do người dùng sinh ra trên các nền tảng mạng xã hội và diễn đàn ẩn chứa những thông tin giá trị về trạng thái tâm lý của chính họ. Tuy nhiên, các bài toán phân loại cảm xúc truyền thống (chỉ đánh giá đơn nhãn hoặc chia ba mức độ: tích cực, tiêu cực, trung tính) chưa đủ sức phản ánh sự phức tạp trong nội tâm con người, khi một câu nói có thể chứa nhiều sắc thái cùng lúc. Vì vậy, đề tài này được thực hiện với hai mục tiêu cốt lõi: (1) nghiên cứu, xây dựng hệ thống phân loại đa nhãn có khả năng nhận diện 28 trạng thái cảm xúc khác nhau của người dùng từ văn bản và (2) triển khai thực nghiệm, đánh giá, so sánh hiệu quả của các thuật toán học máy trong xử lý ngôn ngữ tự nhiên.
 
Để đảm bảo tính khoa học và tập trung, đề tài được giới hạn trong các đối tượng và phạm vi cụ thể. Về mặt dữ liệu, nghiên cứu sử dụng bộ dữ liệu chuẩn GoEmotions do Google cung cấp, tập trung phân tích các bình luận tiếng Anh. Về phương pháp biểu diễn ngôn ngữ, hệ thống tập trung khai thác và so sánh hai kỹ thuật trích xuất đặc trưng văn bản bao gồm phương pháp thống kê truyền thống TF-IDF và kỹ thuật nhúng từ (Word Embedding). Về phương pháp huấn luyện, đề tài tiến hành đánh giá hiệu năng trên tập hợp các thuật toán bao gồm: Logistic Regression, Random Forest, Linear SVC,  Mạng nơ-ron tích chập (CNN) và Mạng bộ nhớ dài - ngắn (LSTM).

Để làm rõ hơn về đề tài này, bài báo cáo được chia làm các nội dung như sau:

\textbf{Chương 1: Tổng quan.}

\textbf{Chương 2: Các cơ sở lý thuyết.}

\textbf{Chương 3: Thực nghiệm.}

\textbf{Chuong 4: Kết quả.}

\textbf{Chương 5: Kết luận.}